% Teoria Micro da Escolha Locacional — PMM-E (Lei 15.233/2025)
% VERSÃO CURTA: 6 slides (capa + 4 conteúdo + takeaway)
% Template: Beamer + Metropolis. Compilar: pdflatex (2x) ou Overleaf.
% Versão longa (10 slides): slides_teoria_micro_beamer.tex
\documentclass[10pt,aspectratio=169]{beamer}

\usetheme[progressbar=frametitle,numbering=fraction]{metropolis}
\usepackage[brazil]{babel}
\usepackage[utf8]{inputenc}
\usepackage[T1]{fontenc}
\usepackage{amsmath,amssymb,mathtools,booktabs,graphicx,hyperref}
\hypersetup{colorlinks=true,linkcolor=violet,urlcolor=violet}

\title{Por que médicos escolhem onde atuar?}
\subtitle{Teoria micro da escolha locacional \\ PMM-E \,—\, Lei 15.233/2025}
\author{Insper \,—\, Economia da Saúde}
\date{\today}
\institute{Moehling et al.\ (2020) \;\; Choné \& Ma (2011) \;\; Reinhardt (1975)}

\begin{document}

\maketitle

%----------------------------------------------------------
\begin{frame}{O puzzle e o núcleo canônico}
\framesubtitle{Vagas de alto IVS ociosas: que custo latente $\Delta c_0$ a bolsa $\Delta B$ supera?}
\begin{equation*}
\arg\max_{i \in I}\; U(\omega_i)
=
\arg\max_{i \in I}\;
\Biggl\{
\sum_{t} \delta^{t}
\left[
\frac{\mathbb{E}\!\left(w_{it}^{(s)}\right)}{p_{it}}
- c_{it}^{(s)}
\right]
\Biggr\}
\quad\text{\footnotesize (Moehling et al.\ 2020, eq.\,1)}
\end{equation*}
\begin{columns}[T]
\begin{column}{0.5\textwidth}
{\small
\textbf{PMM-E:} $V_{im} = \sum_t \delta^t\!\left[\frac{\mathbb{E}(w \mid B_m(\text{IVS}_m))}{p} - c_0(\text{IVS}_m)\right] + \varepsilon_{im}$; bolsa R\$ 10k/15k/20k compõe $w$. Preencher exige $\Delta B / p > \Delta c_0$.
}
\end{column}
\begin{column}{0.5\textwidth}
{\small
\textbf{Subalocação:} capital $w = B + w^{\text{priv}}$ (complementa em 20h) \emph{vs.}\ interior $w = B$ (colapsa no piso) $\Rightarrow$ $B_{\text{interior}} > B_{\text{capital}}$ pode dar $w_{\text{total}}$ menor.
}
\end{column}
\end{columns}
\end{frame}

%----------------------------------------------------------
\begin{frame}{Custos: geográfico e laboral}
\begin{columns}[T]
\begin{column}{0.5\textwidth}
\textbf{Geográfico} {\footnotesize (RRH 2017, eq.\,24)}
\begin{equation*}
c^{\text{geo}}_{im} = \phi(\text{dist}_{im}) - \gamma A_m + \theta_i^{\text{rural}}
\end{equation*}
{\small
$\phi' > 0$ (distância à família dói), $-\gamma < 0$ (amenidade atrai). Limite CNES: sem residência $\Rightarrow$ município do estabelecimento.
}
\end{column}
\begin{column}{0.5\textwidth}
\textbf{Laboral em U} {\footnotesize (Choné \& Ma 2011, eq.\,1)}
\begin{equation*}
c(q) = C(q) - \alpha B(q), \quad c'' \gg 0
\end{equation*}
{\small
$C'>0, C''>0$ (cansaço); $B'>0, B''<0$; $\alpha \ge 0$ (altruísmo). Mínimo em $q_{c_{\max}}$ ($C' = \alpha B'$), neutro em $q_{c=0}$ ($C = \alpha B$).
}
\includegraphics[width=0.85\linewidth]{figuras/curva_custo_laboral_burnout.png}
\end{column}
\end{columns}
\end{frame}

%----------------------------------------------------------
\begin{frame}{$K$, $L$ e a equação-mestra}
\begin{columns}[T]
\begin{column}{0.32\textwidth}
{\small
\textbf{Efeito duplo} (Reinhardt 1975):
\begin{equation*}
\frac{\partial c}{\partial K} < 0, \quad \frac{\partial c}{\partial L} < 0
\end{equation*}
$K$ (centro cirúrgico, UTI) e $L$ (equipe) reduzem $C$ e multiplicam $B$. Sem $K$: resolutividade $\to 0$. Cirúrgico exige $K$; clínico tolera posto simples.
}
\end{column}
\begin{column}{0.68\textwidth}
{\footnotesize
\begin{equation*}
\arg\max_{m}\; U_i = \arg\max_{m}\; \sum_t \delta^t \left[ \frac{\mathbb{E}(w_{mt}^{(s)})}{p_{mt}} - c_{im}^{(s)} \right]
\end{equation*}
\begin{equation*}
\begin{aligned}
c_{im}^{(s)} &=
\underbrace{\phi(\text{dist}_{im}) - \gamma A_m + \theta_i^{\text{rural}}}_{\text{geográfico}} \\
&\quad +
\underbrace{C(q_{im}; L_m, K_m, s) - \alpha_i B(q_{im}; L_m, K_m)}_{\text{laboral líquido}}
\end{aligned}
\end{equation*}
}
\end{column}
\end{columns}
\end{frame}

%----------------------------------------------------------
\begin{frame}{Do latente ao observável — e 4 hipóteses}
\begin{columns}[T]
\begin{column}{0.42\textwidth}
{\small
$c_m = c_0(\text{IVS}_m)$ (IVS 2010, IPEA — 100\% dos municípios):
\textbf{Infra} $\Rightarrow c \uparrow$; \textbf{Cap.\,Humano} $\Rightarrow B' \uparrow$ mas $K \downarrow$; \textbf{Renda/Trab.} $\Rightarrow w = B$. Forças opostas: \textbf{não} assumir $c_0'(\text{IVS}) > 0$.
}
\end{column}
\begin{column}{0.58\textwidth}
{\small
\begin{enumerate}
  \itemsep0pt
  \item \textbf{H1 — Compensação:} $\Delta B/p > \Delta c_0$? Vacância na Faixa-1 (20k) $\Rightarrow$ custo $> $ R\$ 5k.
  \item \textbf{H2 — Subalocação:} isolado ($w = B$) $<$ polo ($w > B$), CNES 6/12m.
  \item \textbf{H3 — Cir.\ $\times$ clín.:} $\left.\frac{\partial \Pr}{\partial \text{IVS}}\right|_{\text{cir.}} \ll \left.\frac{\partial \Pr}{\partial \text{IVS}}\right|_{\text{clín.}}$
  \item \textbf{H4 — IVS:} infra $\downarrow$ fixação; cap.\,hum.\ atenuado/positivo.
\end{enumerate}
}
\end{column}
\end{columns}
\end{frame}

%----------------------------------------------------------
\begin{frame}{Takeaway}
\begin{block}{Mensagem central}
A bolsa compensa, mas não resolve $K$, $L$ e isolamento.
\end{block}
\begin{itemize}
  \item Preencher vaga exige $\Delta B / p > \Delta c_0$ — testar na descontinuidade.
  \item Interior sem mercado privado e sem $K$ afasta cirurgiões primeiro.
  \item Decompor o IVS evita anulação de efeitos opostos.
\end{itemize}
\vspace{0.5em}
{\footnotesize
\textbf{Refs:} Moehling et al.\ (2020); Choné \& Ma (2011); Redding \& Rossi-Hansberg (2017); Reinhardt (1975); IPEA (2015); Costa--Nunes--Sanches (2024).
}
\end{frame}

\end{document}
